% PAGES: 289-289
% Straightforward continuation; no environments open or close on this page.
% "asscts" (below, "for $n=2$ asscts") is printed exactly that way in the
% source (missing "e") -- transcribed as printed, not corrected.
% "(10.7)" here is the same 2x2-matrix-inverse cross-reference already
% flagged as genuine in sec09_c2_2.tex.

Since $\Sigma_R$ is a symmetric positive definite matrix, it follows that
$D^2F_0(w_0)$ is symmetric positive definite. From Theorem 10.7, we find that $w_0$
given by (9.75) is a constrained minimum point.

We conclude that the portfolio with asset allocation $w_m = w_0$, i.e., with
\begin{equation}
w_m \;=\; \frac{1}{\bfone^t\Sigma_R^{-1}\bfone}\Sigma_R^{-1}\bfone,
\end{equation}
is the minimum variance portfolio. Its variance is
\begin{align}
\sigma_m^2 \;=\; w_m^t\Sigma_Rw_m &=
\frac{1}{\left(\bfone^t\Sigma_R^{-1}\bfone\right)^2}\left(\Sigma_R^{-1}\bfone\right)^t
\Sigma_R\left(\Sigma_R^{-1}\bfone\right) \notag\\
&= \frac{1}{\left(\bfone^t\Sigma_R^{-1}\bfone\right)^2}
\bfone^t\left(\Sigma_R^{-1}\right)^t\left(\Sigma_R\Sigma_R^{-1}\right)\bfone \;=\;
\frac{1}{\left(\bfone^t\Sigma_R^{-1}\bfone\right)^2}\bfone^t\Sigma_R^{-1}\bfone
\notag\\
&= \frac{1}{\bfone^t\Sigma_R^{-1}\bfone};
\end{align}
here, we used the fact that $\left(\Sigma_R^{-1}\right)^t = \Sigma_R^{-1}$, since
$\Sigma_R$ is a symmetric matrix; see (9.50).

We conclude this section by deriving the asset allocation for the minimum variance
portfolio made of two assets by using formula (9.76) for $n=2$ asscts.

For $n=2$, the covariance matrix $\Sigma_R$ of the returns of two assets is given by
\[
\Sigma_R \;=\; \begin{pmatrix} \sigma_1^2 & \sigma_1\sigma_2\rho_{1,2} \\
\sigma_1\sigma_2\rho_{1,2} & \sigma_2^2 \end{pmatrix};
\]
see (7.12). From (10.7), it follows that
\begin{equation}
\Sigma_R^{-1} \;=\; \frac{1}{\sigma_1^2\sigma_2^2(1-\rho_{1,2}^2)}
\begin{pmatrix} \sigma_2^2 & -\sigma_1\sigma_2\rho_{1,2} \\ -\sigma_1\sigma_2\rho_{1,2}
& \sigma_1^2 \end{pmatrix}.
\end{equation}

Let $\bfone = \begin{pmatrix} 1 \\ 1 \end{pmatrix}$. Then,
\[
\Sigma_R^{-1}\bfone \;=\; \frac{1}{\sigma_1^2\sigma_2^2(1-\rho_{1,2}^2)}
\begin{pmatrix} \sigma_2^2-\sigma_1\sigma_2\rho_{1,2} \\
\sigma_1^2-\sigma_1\sigma_2\rho_{1,2} \end{pmatrix},
\]
and therefore we find that
\begin{equation}
\bfone^t\Sigma_R^{-1}\bfone \;=\; \frac{\sigma_1^2-2\sigma_1\sigma_2\rho_{1,2}+
\sigma_2^2}{\sigma_1^2\sigma_2^2(1-\rho_{1,2}^2)}.
\end{equation}

From (9.76) for $n=2$, and using (9.78) and (9.79), we obtain that
\begin{align*}
w_m \;=\; \begin{pmatrix} w_1 \\ w_2 \end{pmatrix} &=
\frac{1}{\bfone^t\Sigma_R^{-1}\bfone}\Sigma_R^{-1}\bfone \\
&= \frac{1}{\sigma_1^2-2\sigma_1\sigma_2\rho_{1,2}+\sigma_2^2}
\begin{pmatrix} \sigma_2^2-\sigma_1\sigma_2\rho_{1,2} \\
\sigma_1^2-\sigma_1\sigma_2\rho_{1,2} \end{pmatrix} \\
&= \begin{pmatrix}
\dfrac{\sigma_2^2-\sigma_1\sigma_2\rho_{1,2}}{\sigma_1^2-2\sigma_1\sigma_2\rho_{1,2}+
\sigma_2^2} \\[2ex]
\dfrac{\sigma_1^2-\sigma_1\sigma_2\rho_{1,2}}{\sigma_1^2-2\sigma_1\sigma_2\rho_{1,2}+
\sigma_2^2}
\end{pmatrix}.
\end{align*}
